% JAC1964-20261121. /root/programme_audit; GPT-6 Astra Pro.
% New deduction from the actual projector P13/P17 and curve parameters.
\subsection*{JAC-P18. A fourfold obstruction to an algebraic Fermat radical}

\begin{proposition}[JAC-P18]\label{jac:P18}
Let $F$ be a number field over which the actual Fermat summand
$N=(J,\pi_N)$ of JAC-P13 and JAC-P17 is defined. In particular one
may take $F=\mathbf Q$ or $\mathbf Q(i)$. Put
\[
 \mathcal A_F=\{z\in\operatorname{CH}^2(N_F)_{\mathbf Q}:
       z_{\overline F}\text{ is algebraically trivial on }
       J_{\overline F}\}.
\]
If $\mathcal A_F\ne0$, there are a finite extension $F'/F$, a
smooth projective geometrically integral curve $C/F'$ with a point
$c_0\in C(F')$, and a nonzero cycle
\[
 Z\in\operatorname{CH}^2(C\times J_{F'})_{\mathbf Q}
\]
such that
\[
  Z=(\pi_1^C\otimes\pi_N)_*Z,
       \qquad \operatorname{cl}(Z)=0.
                                                        \tag{P18.1}
\]
Here
$\pi_1^C=\Delta_C-[c_0]\times C-C\times[c_0]$,
and homological triviality holds for Betti cohomology under every
complex embedding and for every $\ell$-adic realization.
The cycle $Z$ belongs to the degree-four component of the product
Chow--K\"unneth decomposition on the fourfold $C\times J_{F'}$:
\[
                 (\pi_4^{C\times J})_*Z=Z.             \tag{P18.2}
\]

Consequently, injectivity of the ordinary cycle-class map on
$(\pi_1^C\otimes\pi_N)_*\operatorname{CH}^2(C\times J_{F'})_{
\mathbf Q}$ for every such $F'$ and $C$ implies
$\mathcal A_F=0$. Murre's conjecture~(D) in codimension two for
these product decompositions is a sufficient condition. No
converse and no injectivity on the full group
$\operatorname{CH}^2(N_F)_{\mathbf Q}$ is asserted.
\end{proposition}

\begin{proof}
Take $0\ne z\in\mathcal A_F$. Algebraic equivalence is generated
by differences of fibres of cycle families over smooth projective
curves. Thus over $\overline F$ there is an equality
\[
 z_{\overline F}=\sum_j(\Gamma_j)_*D_j,
 \qquad
 \Gamma_j\in\operatorname{CH}^2(C_j\times J)_{\mathbf Q},
 \quad D_j\in\operatorname{Pic}^0(C_j)_{\mathbf Q}.
                                                        \tag{P18.3}
\]
One may take each $D_j$ to be a rational linear combination of
differences of rational points. There are only finitely many
curves, cycles, points and rational equivalences in this
expression. Enlarge $F$ by a finite extension $F'$ so that all
are defined over $F'$, all the curves are smooth projective and
geometrically integral, and a chosen basepoint on each curve is
$F'$-rational. The equality then holds over $F'$ after enlarging
$F'$ once more if necessary. This last assertion follows directly
from the finite data in a rational equivalence, or from the
direct-limit description of Chow groups under algebraic field
extension.

Rational restriction from $F$ to $F'$ is injective, because its
composition with norm is multiplication by $[F':F]$.
Hence $z_{F'}\ne0$. Apply $\pi_N$ to (P18.3). Since
$(\pi_N)_*z_{F'}=z_{F'}$, at least one summand
$\pi_N\circ\Gamma_j$ has nonzero action on $D_j$. Choose that
summand, and write its curve, divisor and correspondence as
$C,D,\Gamma$.

The projector $\pi_1^C$ acts identically on
$\operatorname{Pic}^0(C)_{\mathbf Q}$. Define the actual Chow
correspondence
\[
              Z=\pi_N\circ\Gamma\circ\pi_1^C.
                                                        \tag{P18.4}
\]
It still has codimension two on $C\times J$: a correspondence
from the curve to the threefold that takes divisors to
codimension-two cycles has precisely this codimension.
Its action on $D$ is the chosen nonzero summand, so $Z\ne0$ as
a rational Chow class. This uses its actual action, not its
numerical or cohomological realization.

When composition of correspondences is viewed as an action on
the product cycle, the source projector is transposed:
\[
 \pi_N\circ\Gamma\circ\pi_1^C
       =({}^t\pi_1^C\otimes\pi_N)_*\Gamma.
\]
For the displayed degree-one curve projector
${}^t\pi_1^C=\pi_1^C$. Idempotence of both projectors proves
the first identity in (P18.1).

Fix an embedding $F'\hookrightarrow\mathbf C$. The cohomology
class of $Z$ therefore belongs to
\[
 H^1(C,\mathbf Q)\otimes H^3(N,\mathbf Q)
       \ \subset\ H^4(C\times J,\mathbf Q).
\]
The two factors have Hodge types $(1,0),(0,1)$ and
$(3,0),(0,3)$, respectively. Their tensor product has only types
$(4,0),(3,1),(1,3),(0,4)$ and has no $(2,2)$ part. An algebraic
codimension-two cycle has type $(2,2)$, so
$\operatorname{cl}(Z)=0$. The same argument works for every
embedding. Compatibility of Betti and $\ell$-adic cycle classes
gives the asserted vanishing in every $\ell$-adic realization.

The product decomposition has
\[
 \pi_4^{C\times J}
   =\pi_0^C\otimes\pi_4^J
      +\pi_1^C\otimes\pi_3^J
      +\pi_2^C\otimes\pi_2^J.
\]
JAC-P13/P17 gives
$\pi_N\pi_3^J=\pi_3^J\pi_N=\pi_N$ in rational Chow
correspondences. Orthogonality now proves (P18.2).
Murre~(D) in this codimension says that
$\ker(\pi_4^{C\times J})_*$ is exactly the homologically
trivial subgroup. It would force the nonzero $Z$ simultaneously
to be killed and fixed by $\pi_4^{C\times J}$, a contradiction.
The weaker stated injectivity on the single component gives the
same contradiction and proves the sufficient condition.
\end{proof}

\paragraph{Source and logical scope.}
The actual projector and the extreme Hodge types are the proved
JAC-P13/P17 inputs. Curve parameterization of algebraic
equivalence is used in its usual form, as in Achter--Casalaina-Martin--Vial,
\emph{Parameter spaces for algebraic equivalence},
Proposition~3.5 and Lemma~3.9, author pp.~4--7,
\url{https://www.math.uni-bielefeld.de/~vial/abtriv.pdf}.
The required product assertion is not supplied by the theorem of
Xu--Xu: \emph{Remarks on Murre's conjecture on Chow groups},
Corollary~3.9, author p.~9, proves (A) and (B) for an abelian
threefold times a curve;
\url{https://arxiv.org/pdf/1110.0668}.
Kenichiro Kimura, \emph{Murre's conjectures for certain product
varieties}, author p.~2 and Theorem~2.1, p.~3, explicitly isolates
the related degree-four component and treats restricted support
types under additional hypotheses;
\url{https://arxiv.org/pdf/math/0509249}.
Those hypotheses do not cover $C\times E_i^3$.
The obstruction here concerns rational Chow faithfulness in one
degree-four component; it must not be confused with the standard
conjecture identifying homological and numerical equivalence.
It refines AR-15a by keeping the obstructing variety four-dimensional.
